	% An Overview of Lie Algebra's
% Adapted from Justin Ryan's "Examples of Lie Algebras"
% Chaskin Saroff
% last changed: 29 April 2015
% feel free to make any improvements/changes you wish

\documentclass[9 pt]{beamer}
\usepackage{etex}
\reserveinserts{28}
% \usepackage{default}
\usepackage{lmodern}
\usepackage{amsmath,amsfonts,epsfig,pgf} % ,graphicx
\usepackage{graphicx}

% choose your theme
 \usetheme{Warsaw} % Warsaw, Copenhagen, Darmstadt, Madrid, Singapore, etc...

% custom SUNY Oswego color scheme
\definecolor{oswego}{rgb}{0.15,0.4,0.15}
\setbeamercolor*{palette primary}{fg=white, bg=oswego}
\setbeamercolor*{palette sidebar primary}{fg=black, bg=oswego}
\setbeamercolor{block title}{bg=black,fg=white} % bg=background, fg= foreground
\setbeamercolor{block body}{bg=oswego,fg=black} % bg=background, fg= foreground
\setbeamercolor{alerted text}{fg=oswego}
\usecolortheme[named={oswego}]{structure}
\def\today{\number\day\space\ifcase\month\or
   January\or February\or March\or April\or May\or June\or
   July\or August\or September\or October\or November\or December\fi
   \space\number\year}

% something I found to get alert blocks in the Oswego color scheme
\newenvironment<>{lakeblock}[1]{%
  \begin{actionenv}#2%
      \def\insertblocktitle{#1}%
      \par%
      \mode<presentation>{%
\setbeamercolor{block title}{fg=white,bg=black}
       \setbeamercolor{block body}{fg=white,bg=oswego}
            }%
      \usebeamertemplate{block begin}}
    {\par\usebeamertemplate{block end}\end{actionenv}}

% Oswego State logo on every page
%\logo{\pgfputat{\pgfxy(-10.5,0.35)}{\pgfbox[center,center]
%{\includegraphics[height=1.35cm]{oswego_logo_horiz_grn}}}}

% commutative diagrams with XY-pic
\usepackage[all]{xy}
\SelectTips{cm}{}
% make \mathscr, TeX \cal, and Euler script *all* available
% (notice the new command names to avoid overlap and/or confusion)
\usepackage{mathrsfs}
\let\rscr=\mathscr % use \rscr{} for Ralph Smith fancy script
\let\mathscr=\relax
\let\mcal=\mathcal % use \mcal{} for TeX \cal script
\usepackage{eucal}
\let\escr=\mathcal % use \escr{} for Euler script
\let\mathcal=\relax
% a better "bar" thanks to Donald Arsenau -- see \pbar infra
\usepackage{accents}

% title page information
\title[Encoding the Bracket - An Overview of Lie Agebra's]{Encoding the Bracket - An Overview of Lie Algebras}
\author[A.\, Jansing \,C.\, Saroff]{Alexander Jansing,  Chaskin Saroff}
\institute[SUNY Oswego]{Oswego State University\\ Department of Mathematics}
\date{\today}

% new math commands
\newcommand{\at}[1]{\emph{\alert{#1}}}
\newcommand{\ad}[1]{\text{ad}_{#1}}
\newcommand{\add}[1]{\ad{#1}^\dagger}
\newcommand{\br}[2]{\left[ #1, #2 \right]}
\newcommand{\bre}{\br{\ }{\,}}
\newcommand{\C}{\mathbb{C}}
\newcommand{\F}{\mathbb{F}}
\newcommand{\h}{\lag{h}}
\newcommand{\inp}[2]{\langle #1, #2 \rangle}
\newcommand{\inpe}{\inp{\ }{\,}}
\newcommand{\lag}[1]{\mfrak{#1}}
\newcommand{\mfrak}[1]{\mathfrak{#1}}
\newcommand{\R}{\mathbb{R}}
\renewcommand{\a}{\alpha}
\newcommand{\surj}{\rightarrow\kern-.82em\rightarrow}
\newcommand{\tQ}{\widetilde{Q}}
\renewcommand{\v}{\lal{v}}
\newcommand{\z}{\lal{z}}
\newcommand{\V}{\mathfrak{g}}
\newcommand{\fg}{\mathfrak{g}}
\newcommand{\fz}{\mathfrak{z}}
\newcommand{\fv}{\mathfrak{v}}
\newcommand{\fh}{\mathfrak{h}}
\newcommand{\QQ}{\mathbb{Q}}
\newcommand{\ZZ}{\mathbb{Z}}
\newcommand{\RR}{\mathbb{R}}
\newcommand{\CC}{\mathbb{C}}
\newcommand{\NN}{\mathbb{N}}
\newcommand{\FF}{\mathbb{F}}
\newcommand{\zvec}{\mathbf{0}}
\newcommand{\lal}[1]{\mathfrak{#1}}
\newcommand{\lan}{\lal{n}}
\newcommand{\lav}{\lal{v}}
\newcommand{\laz}{\lal{z}}


% colored text commands
\newcommand{\red}[1]{{\color{red} #1}}
\newcommand{\grn}[1]{{\color{green} #1}}
\newcommand{\blu}[1]{{\color{blue} #1}}
\newcommand{\ylw}[1]{{\color{yellow} #1}}
\newcommand{\mgn}[1]{{\color{magenta} #1}}
\newcommand{\cyn}[1]{{\color{cyan} #1}}

%Stuff from Justin's example
%\renewcommand{\a}{\vect{a}}
%\newcommand{\abs}[1]{\vert #1 \vert}
%\renewcommand{\b}{\vect{b}}
%\newcommand{\bve}{\pbar{\ve}}
%\newcommand{\dps}{\displaystyle}
%\newcommand{\ds}{\oplus}
%\newcommand{\name}{{\bf Name:}\underline{\hspace{3 in}}}
%\newcommand{\nin}{\noindent}
%\newcommand{\norm}[1]{\left\Vert #1 \right\Vert}
%\newcommand{\pg}[1]{\paragraph{#1}}
%\newcommand{\ph}[1]{\phantom{#1}}
%\renewcommand{\span}[1]{\textrm{span}\left\{ #1 \right\}}
\renewcommand{\t}{\texttt{t}}
\newcommand{\T}{\textsc{T}}
\renewcommand{\u}{\vect{u}}
\renewcommand{\v}{\vect{v}}
\newcommand{\ve}{\varepsilon}
\newcommand{\vc}[1]{\left\langle #1 \right\rangle}
\newcommand{\vect}[1]{\mathbf{#1}}

% tiks packages
\usepackage{tikz}
\usetikzlibrary{calc}
\usepackage{tikzscale}
\usepackage{filecontents}
\usepackage{caption}
\usepackage{subcaption}
\captionsetup{compatibility=false}
\usepackage{wrapfig}

\begin{filecontents*}{heis.tikz}
    \begin{tikzpicture}
        \draw (-4,4) -- (1,2);
        \draw (1, 2) -- (4,4);
        \draw (-1,6) -- (4,4);
        \draw (-1,6) -- (-4,4);
        \draw [-] (0,8) -- (0,4);
        \draw [-] (0,2.4) -- (0,0);
        %\draw [dashed] (4,0) -- (4,3);
        \node [right] at (0, 8) {$\fz$};
        \node [below] at (4,4) {$\fv$};
    \end{tikzpicture}
\end{filecontents*}

\newcommand{\tikzmark}[1]{\tikz[overlay,remember picture] \node (#1) {};}
\newcommand{\DrawBox}[1][]{%
    \tikz[overlay,remember picture]{
    \draw[red,#1]
      ($(left)+(-0.2em,0.9em)$) rectangle
      ($(right)+(0.2em,-0.3em)$);}
}

\NewDocumentCommand{\highlight}{O{blue!40} m m}{%
    \draw[mycolor=#1] (#2.north west)rectangle (#3.south east);
}

\NewDocumentCommand{\fhighlight}{O{blue!40} m m}{%
    \draw[myfillcolor=#1] (#2.north west)rectangle (#3.south east);
}

\begin{document}

\section{Definition and Examples}


\begin{frame}{}
\vspace{1.35 cm}

 \titlepage

\includegraphics[height=1.35cm]{oswego_logo_horiz_grn}

\end{frame}

% \begin{frame}{}
%
% \end{frame}

% \begin{frame}{Outline}
%   \begin{center}
%     \begin{minipage}[]{.85\textwidth}
%       \tableofcontents
%     \end{minipage}
%   \end{center}
% \end{frame}

\subsection{Definition}

\note{So the Lie Algebra is a \emph{vector space}. Does anyone recognize
a similar operator, maybe from calulus or physics, that takes two vectors
and produces a vector?}
\begin{frame}{Definition}
A \emph{\alert{Lie Algebra}} is a (finite-dimensional) vector space $\fg$
together with a bilinear multiplication $\bre : \fg \times \fg \to \fg$
that satisfies three properties: \pause

\smallskip

\begin{enumerate}
\item $\br{x}{x} = 0$ for all $x \in \fg$,
\pause

\

\item $\br{\br{x}{y}}{z} + \br{\br{y}{z}}{x} + \br{\br{z}{x}}{y} = 0$ for
all $x, y, z \in \fg$.
and
\pause

\

\item  $\br{x+y}{z} = \br{x}{z}+\br{y}{z}\quad \text{and} \quad 
    \br{\a x}{y}=\br{x}{\a    y}=\a\br{x}{y}
    \ \forall x,y,z \in \fg \ \text{and} \ \a \in \RR$ 
\end{enumerate}
\pause

\

Property 1 is known as \emph{\alert{alternating}}, property 2 is known
as the \emph{\alert{Jacobi Identity}} and property 3 describes
\emph{\alert{bilinearity}}.
\pause

\

Property 1 and 3 imply another property called \emph{\alert{skew symmetry}}\,:

\smallskip

\begin{enumerate}
\item $\br{x}{y} = -\br{y}{x}$ for all $x, y \in \fg$. 
\end{enumerate}
\pause 

\

The product $\bre$ is known as a \emph{\alert{Lie bracket}} on $\fg$.

\end{frame}

\subsection{Examples}

\begin{frame}{skew-symmetry}
By bilinearity, every alternating product is also skew-symmetric.
Indeed, if [ , ] is alternating then
\pause
        \begin{align}
            0 &= [x + y, x + y] 
            \\&= [x + y, x] + [x + y, y]
            \\&= [x,x] + [y, x] + [x + y, y]
            \\&= [x,x] + [y, x] + [x, y] + [y, y]
            \\0&=         [x, y] + [y, x]
            \\\implies [x,y] &= -[y,x]
        \end{align}
\pause
Conversely, if $\bre$ is skew-symmetric, then
$\br{x}{x}+\br{x}{x}$ = 0 implies that 2[x,x] = 0. Now we see that this implies [x,x] = 0
so long as our field is not of characteristic 2, for in those spaces 2 = 0 and we
can deduce nothing about [x,x].
\end{frame}

\begin{frame}{The Center}
    %\begin{definition}
        The \emph{center} of a Lie Algebra, $\fg$ is
        \pause
        \\\[\fz=\{z \in \fg \mid [z,x]
        = \zvec \; \forall x \in \fg\}.\]
        \pause
        \\Thus $z$, is in the center of $\fg$ if and only if
        \\\[[z,x] = \zvec \; \forall x \in \fg.\]

        The \emph{non-center}, $\fv$, is given by $\fv = \fg - \fz$
    %\end{definition}
\end{frame}

\begin{frame}{A First Example: One-Step Nilpotent Lie Algebras}
Let $\fg$ be any vector space, and define
$$
\br{x}{y} = 0 \text{\ \ for\ all\ } x,y \in \fg.
$$
\pause

\

Clearly $\bre$ is alternating:
$$
\br{x}{x} = 0 \text{\ \ for\ all\ } x \in \fg.
$$
\pause

\

And the Jacobi Identity is trivial:

\pause
\begin{eqnarray*}
\br{\br{x}{y}}{z} + \br{\br{y}{z}}{x} + \br{\br{z}{x}}{y} & = & 
	\br{0}{z} + \br{0}{x} + \br{0}{y} \\
	& = & 0 + 0 + 0 = 0.
\end{eqnarray*}
\pause
\
This Algebra is called \emph{Abelian} or \emph{One-Step Nilpotent}.
\pause

\end{frame}

\begin{frame}{Two-Step Nilpotent}

    A Lie Algebra, $\fg$ is called \emph{two-step nilpotent} if and only if
    \[
        [[x,y],z] = \zvec \; \forall x,y,z \in \fg
        \ \text{and} \ \fz \neq \fg
    \]
    \pause
    \begin{enumerate}
        \item Applying the bracket twice to any arbitrary set of three vectors in the Algebra, always returns $\zvec$.
        \pause
        \item Applying the bracket once to any arbitrary set of two vectors does not always return $\zvec$.
    \end{enumerate}
\end{frame}

\begin{frame}{$\R^3$ as a non-Abelian Lie Algebra}
Let $x = (x_1, x_2, x_3)$ and $y = (y_1, y_2, y_3)$. The \at{cross product}
of $x$ with $y$ is defined by
$$
x \times y = (x_2y_3 - x_3y_2, x_3y_1 - x_1y_3, x_1y_2 - x_2y_1).
$$
\pause

The cross product is skew symmetric:
$$
x \times y = - y \times x \text{\ \ for\ all\ } x,y \in \R^3.
$$
\pause

Using some identities from Calc III, we can show that the cross product satisfies
the Jacobi Identity:
\pause
\begin{eqnarray*}
& & \hspace{-1.5 cm} (x \times y) \times z + (y \times z) \times x + (z \times x) \times y
\pause
\\ & = & 
\red{(x \cdot z)y} - \blu{(y \cdot z)x} + \grn{(y\cdot x)z} - \red{(z \cdot x)y} +
\blu{(z \cdot y)x} - \grn{(x \cdot y)z}
\pause
\\ & = & 0.
\end{eqnarray*}
\end{frame}

\begin{frame}{$\h_3$ : The Heisenberg Algebra}
We can construct another example from $\RR^3$ by defining a different Lie Bracket.
Consider a 3-dimensional vector space with basis vectors $x, y, z$.
\pause
Define a Lie bracket on this space by
$$
\br{x}{y} = z,
$$
with all other brackets equal to 0.
\pause
This is called the \at{Heisenberg Algebra}, and is denoted by $\lag{h}_3$.
\pause

\begin{figure}[h]
    \centering
    \includegraphics[width=30mm]{heis.tikz} %input file
    \label{fig:pic} %label name
    \caption{\footnotesize The Heisenberg Algebra, $\fh_3$} %caption
\end{figure}

\pause

Can anyone tell me if the Heisenberg Algebra is one-step nilpotent,
two-step nilpotent, neither, or both?

\end{frame}

\section{Defining the Lie Algebra}

\subsection{Inner Products}

\begin{frame}{Inner Products}
An \at{inner product} on a vector space $\fg$ is a non-degenerate, symmetric,
bilinear, postive definite function $\inpe :\fg \times \fg \to \R$.
\pause
\begin{enumerate}
\item If $\inp{x}{y} = 0$ for all $y \in \fg$, then $x$ must be $0$;
\pause

\item $\inp{x}{y} = \inp{y}{x}$ for all $x,y \in \fg$;
\pause

\item $\inp{ax + by}{z} = a\inp{x}{z} + b\inp{y}{z}$ for all $x,y,z \in \fg$, $a, b, \in \R$.
\item $\inp{x}{x} \ge 0 \ \forall x \in \fg$ and $\inp{x}{x}=0 \implies x=\zvec$
\end{enumerate}
\pause

A Lie Algebra with an inner product is called a \emph{Metric Lie Algebra}
\pause

\

{\bf Example} The \at{dot product} on $\R^3$:
\pause
Let $x, y \in \R^3$; then $x \cdot y = x^Ty$.  The dot product is given by
\pause
\[
    \langle x,y \rangle =  x_1y_1 + x_2y_2 + x_3y_3
\]

\

\end{frame}

\subsection{Encoding Geometry via Algebra}

\begin{frame}{$j$-maps}
    In order to seperate the interaction between the center and the non-center
    of a two-step nilpotent lie algebra, we can define a map, $j_z:\fv \to \fv$.
    \pause
    \\$j_z$ is defined by the equation
    $$
    \inp{\br{x}{y}}{z} = \inp{y}{j_zx}
    $$
    \pause

    \

    The $j$-maps are extremely useful because they simultaneously encode information
    about the algebraic structure and geometric structure of the Lie algebra $\lag{g}$.
    \pause

    \

    This makes the computation of many seemingly complicated objects
    ``simple" calculations in linear algebra!
\end{frame}

\begin{frame}{The Lie Bracket as a Matrix with Vector entries}
    In order to do math with Lie Algebras on a computer,
    we would like to represent the Lie Bracket in a matrix form.
    \pause
    Consider the Lie Bracket, with bases $e_1, e_2, \dots, e_n$
    \\Formally, $L_{ij} = [e_i, e_j]$
    \pause

    $$
     L = \begin{pmatrix}
        L_{11} & L_{12} & \cdots & L_{1n} \\
        L_{21} & L_{22} & \cdots & L_{2n} \\
        \vdots &        & \ddots & \vdots \\
        L_{n1} & L_{n2} & \cdots & L_{nn}
    \end{pmatrix}
    $$
    \pause
    We will denote the $k$th entry of $L_{ij}$ as $L_{ij}^k$.
\pause
\\$[e_i,e_i] = \zvec$
\end{frame}
\begin{frame}{The Lie Bracket as a Matrix with Vector entries}
Recall that the Lie Bracket is a matrix of vectors.\\
\pause
\begin{figure}[]
    \centering
    \includegraphics[width=50mm]{Lie_MX_Flat.tikz} %input file
    \label{fig:pic} %label name
    \caption{\footnotesize The Lie Bracket as a Matrix, $L$} %caption
\end{figure}

\end{frame}
\begin{frame}{The Lie Bracket as a Matrix with Vector entries}
Recall that the Lie Bracket is a matrix of vectors.\\
\begin{figure}[]
    \centering
    \includegraphics[width=50mm]{Lie_MX_Flat_ij.tikz} %input file
    \label{fig:pic} %label name
    \caption{\footnotesize The Lie Bracket as a Matrix, $L$} %caption
\end{figure}

\end{frame}
\begin{frame}{The Lie Bracket as a Matrix with Vector entries}
Recall that the Lie Bracket is a matrix of vectors.\\
\begin{figure}[]
    \centering
    \includegraphics[width=50mm]{Lie_MX_Vector.tikz} %input file
    \label{fig:pic} %label name
    \caption{\footnotesize The Lie Bracket as a Matrix, $L$} %caption
\end{figure}
\end{frame}
\begin{frame}{The Lie Bracket as a vector of matrices}
    It can also be represented by a ``stack'' of matrices
    \setlength{\unitlength}{0.25 cm}
    \begin{center}
        \begin{picture}(0,0)
            \put(-8,0){\line(1,0){16}}
            \put(-8,0){\line(1,-4){1}}
            \put(-7,-4){\line(1,0){16}}
            \put(8,0){\line(1,-4){1}}
            \put(-8,-2){\line(1,0){0.5}}
            \put(-8,-2){\line(1,-4){1}}
            \put(-7,-6){\line(1,0){16}}
            \put(9,-6){\line(-1,4){0.5}}
            \put(-8,-4){\line(1,0){0.5}}
            \put(-8,-4){\line(1,-4){1}}
            \put(-7,-8){\line(1,0){16}}
            \put(9,-8){\line(-1,4){0.5}}
            \put(10,-4){$L^1$}
            \put(10,-6){$L^2$}
            \put(10,-8){$L^3$}
            \put(10.25,-10){$\vdots$}
            \put(10,-12){$L^m$}
            \put(-7,-10){$\vdots$}
            \put(-8,-7.75){$\vdots$}
            \put(7.75,-10){$\vdots$}
            \put(1,-10){$\vdots$}
            \put(-8,-8){\line(1,0){0.75}}
            \put(-8,-8){\line(1,-4){1}}
            \put(-7,-12){\line(1,0){16}}
            \put(9,-12){\line(-1,4){0.9}}
            \put(-11,-7){$L = $}
        \end{picture}
    \end{center}

    \vspace{3.15 cm}
    A vector of matrices.
    $$
    L_{ij} = \begin{pmatrix}
        L_{ij}^1 \\ L_{ij}^2 \\ \vdots \\ L_{ij}^m
    \end{pmatrix}.
    $$
\end{frame}
\begin{frame}
    $$
    L_{ij} = \begin{pmatrix}
        L_{ij}^1 \\ L_{ij}^2 \\ \vdots \\ L_{ij}^m
    \end{pmatrix}.
    $$
    These $L_{ij}$ can then be arranged in an $n \times n$ matrix
    $$
    L = (L_{ij}) = \begin{pmatrix}
        L_{11} & L_{12} & \cdots & L_{1n} \\
        L_{21} & L_{22} & \cdots & L_{2n} \\
        \vdots & \vdots & \ddots & \vdots \\
        L_{n1} & L_{n2} & \cdots & L_{nn}
    \end{pmatrix}.
    $$

\end{frame}
\begin{frame}
    It's nice to see that we can \emph{encode} the bracket into a matrix
    representation, but how can we use that encoding as the Lie Bracket
    ``map''?
    \pause
    \\Notice that we can retrieve information about the bases.
    \pause
    \\$\br{e_1}{e_2} = L_{12}$
    \pause
    $$
     L = \begin{pmatrix}
         L_{11}& \tikzmark{left}L_{12}\tikzmark{right} & \cdots & L_{1n} \\
            \DrawBox[thick]
            L_{21}& L_{22} & \cdots & L_{2n} \\
               \vdots &       & \ddots & \vdots \\
               L_{n1}&L_{n2}& \cdots & L_{nn}
        \end{pmatrix}
    $$
    \pause
    \\So it will be enough find a way to write all Lie Brackets as a linear
    combination of basis brackets.
\end{frame}

\begin{frame}{Lie Bracket as Linear Combination of Bases}
    %I had to use uncover here instead of pause because align ``does wicked
    %things too beamer''
    \uncover<00->{We can write $\br{x}{y}$ as a linear combination of their component vectors'
brackets.}
    \uncover<01->{\\Let's look at an example of a 2 dimensional space}
    \uncover<02->{\\Let $x,y \in \fg$ with}
    \uncover<03->{\\$x=w_1e_1 + w_2e_2$}
    \uncover<04->{\\$y=v_1e_1 + v_2e_2$}
    \begin{align*}
        \uncover<05->{\text{So} \br{x}{y} &= \br{w_1e_1 + w_2e_2}{v_1e_1 + v_2e_2}}
        \uncover<06->{\\&= \br{w_1e_1}{v_1e_1 + v_2e_2} + \br{w_2e_2}{v_1e_1 + v_2e_2}}
        \uncover<07->{\\&= \br{w_1e_1}{v_1e_1} + \br{w_1e_1}{v_2e_2}+ \br{w_2e_2}{v_1e_1} + \br{w_2e_2}{v_2e_2}}
        \uncover<08->{\\&= w_1v_1\br{e_1}{e_1} + w_1v_2\br{e_1}{e_2}+ w_2v_1\br{e_2}{e_1} + w_2v_2\br{e_2}{e_2}}
        \uncover<09->{\\&= w_1v_1L_{11} + w_1v_2L_{12} + w_2v_1L_{21}+ w_2v_2L_{22}}
    \end{align*}
    \uncover<10->{This means that the Lie Bracket is fully described by the matrix, L}
\end{frame}
\begin{frame}{The Lie Bracket as a Matrix with Vector entries}
    Let's take another look at $L$
    \pause
    \\Consider the Lie Bracket, with bases $e_1, e_2, \dots, e_n$
    \\Formally, $L_{ij} = [e_i, e_j]$
    \pause

    $$
     L = \begin{pmatrix}
        L_{11} & L_{12} & \cdots & L_{1n} \\
        L_{21} & L_{22} & \cdots & L_{2n} \\
        \vdots &        & \ddots & \vdots \\
        L_{n1} & L_{n2} & \cdots & L_{nn}
    \end{pmatrix}
    $$
\pause
What is special about $[e_i,e_i]$?
\pause
\\$[e_i,e_i] = \zvec$
\end{frame}
\begin{frame}{The Lie Bracket as a Matrix with Vector entries}
    Let's take another look at $L$
    \\Consider the Lie Bracket, with bases $e_1, e_2, \dots, e_n$
    \\Formally, $L_{ij} = [e_i, e_j]$
    $$
     L = \begin{pmatrix}
        \zvec& L_{12} & \cdots & L_{1n} \\
        L_{21} & \zvec& \cdots & L_{2n} \\
        \vdots &        & \ddots & \vdots \\
        L_{n1} & L_{n2} & \cdots & \zvec
    \end{pmatrix}
    $$
\pause
Given $[e_i, e_j]$, what do we know about $[e_j, e_i]$?
\pause
\\$[e_i, e_j] = -[e_j, e_i]$
\end{frame}
\begin{frame}{The Lie Bracket as a Matrix with Vector entries}
    Let's take another look at $L$
    \\Consider the Lie Bracket, with bases $e_1, e_2, \dots, e_n$
    \\Formally, $L_{ij} = [e_i, e_j]$
    $$
     L = \begin{pmatrix}
               \zvec& L_{12}  & \cdots & L_{1n} \\
               -L_{12}& \zvec & \cdots & L_{2n} \\
               \vdots &       & \ddots & \vdots \\
               -L_{1n}&-L_{2n}& \cdots & \zvec
        \end{pmatrix}
    $$
    \\
    Consider the map $\br{e_1}{\,}: \fg \to \fg$.
    \\Can anyone see the matrix representation of $\br{e_1}{\,}$?
\end{frame}

\begin{frame}{The Lie Bracket as a Matrix with Vector entries}
    Let's take another look at $L$
    \\Consider the Lie Bracket, with bases $e_1, e_2, \dots, e_n$
    \\Formally, $L_{ij} = [e_i, e_j]$
    $$
     L = \begin{pmatrix}
         \tikzmark{left}\zvec& L_{12}  & \cdots & L_{1n} \tikzmark{right}\\
            \DrawBox[thick]
               -L_{12}& \zvec & \cdots & L_{2n} \\
               \vdots &       & \ddots & \vdots \\
               -L_{1n}&-L_{2n}& \cdots & \zvec
        \end{pmatrix}
    $$
    \\
    Consider the map $\br{e_1}{\,}: \fg \to \fg$.
    \\Can anyone see the matrix representation of $\br{e_1}{\,}$?
\end{frame}
\begin{frame}{Generalizing the Lie Bracket}
    As before, we'll investigate the 2 dimensional example.
    \pause
    Let
    $$
     L = \begin{pmatrix}
            L_{11} & L_{12} \\
            L_{21} & L_{22}
         \end{pmatrix},\ \ 
    $$
    \pause
    with
    $$
     x = \begin{pmatrix}
            w_1\\
            w_2
         \end{pmatrix},\ \ 
     y = \begin{pmatrix}
             v_1\\
             v_2
         \end{pmatrix}
    $$
    \pause
    $$
    x^{T}L = 
    \begin{pmatrix}
        w_1 w_2
    \end{pmatrix}
     \begin{pmatrix}
            L_{11} & L_{12} \\
            L_{21} & L_{22}
     \end{pmatrix}
     =
     \begin{pmatrix}
         w_1L_{11} + w_2L_{12}, w_1L_{12}+w_2L_{22}
     \end{pmatrix}
    $$
    \pause
    $$
        (x^TL)y = 
        \begin{pmatrix}
            w_1L_{11} + w_2L_{12}, w_1L_{12}+w_2L_{22}
        \end{pmatrix}
        \begin{pmatrix}
            v_1 \\
            v_2
        \end{pmatrix}
    $$
    \pause
    $$
    =   \begin{pmatrix}
            w_1v_1L_{11} + w_2v_1L_{12} + w_1v_2L_{12}+w_2v_2L_{22}
        \end{pmatrix}
    $$
    \pause
    That is
    $\br{x}{y} = x^TLy$
\end{frame}
\begin{frame}{Example Abelian $\RR^3$}
    What does Abelian $\RR^3$ look like in this form?
    \pause
    \\Recall that a Lie Algebra is Abelian if
    \pause
    \\$\br{x}{y} = \zvec \ \forall x,y \in \RR^3$
    \pause
    $$
    \begin{pmatrix}
            \zvec & \zvec & \zvec
          \\\zvec & \zvec & \zvec
          \\\zvec & \zvec & \zvec
    \end{pmatrix}
    $$
    \pause
    What is the center of Abelian $\RR^3$
\end{frame}
\begin{frame}{Examples of Cross Product over $\RR^3$}
    What does the the cross product on $\RR^3$ look like in this form?
    \pause
    \\$\br{e_1}{e_2}=e_3$
    \\$\br{e_1}{e_3}=-e_2$
    \\$\br{e_2}{e_3}=e_1$
    \pause
    $$
    \begin{pmatrix}
          \zvec & e_3   & -e_2
        \\-e_3  & \zvec & e_1
        \\e_2   & -e_1  & \zvec
    \end{pmatrix}
    $$
    \pause
    What is the center of $\RR^3$ with the cross product?
\end{frame}
\begin{frame}{$\h_3$'s Lie Bracket}
    How could the three dimensional Heisenberg Algebra be represented?
    \pause
        $\br{e_1}{e_2}=e_3$ for bases $e_1, e_2, e_3 \in \h_3$
        \\And all other brackets zero
    \pause
    $$
    \begin{pmatrix}
          \zvec & e_3   & \zvec
        \\-e_3  & \zvec & \zvec
        \\\zvec & \zvec & \zvec
    \end{pmatrix}
    $$
    \pause
    \\What is the center of $\h_3$?
    \pause
    \\The center is the span of $\{e_3\}$.
\end{frame}
\begin{frame}{Inner Product as a Matrix}
    The Inner Product can be represented by a matrix very similarly to the Lie Bracket.
    \pause
    \\$\inp{e_i}{e_j} = E_{ij}$
    \pause
    $$
        E = \begin{pmatrix}
                E_{11} & E_{12} & \cdots & E_{1n} \\
                E_{21} & E_{22} & \cdots & E_{2n} \\
                \vdots & \vdots & \ddots & \vdots \\
                E_{n1} & E_{n2} & \cdots & E_{nn}
            \end{pmatrix}
    $$
    \pause
    Recall that $\inp{e_i}{e_j} = \inp{e_j}{e_i}$
\end{frame}
%Transition from previous frame showing that e_12 = e_21
\begin{frame}{Inner Product as a Matrix}
    The Inner Product can be represented by a matrix very similarly to the Lie Bracket.
    \\$\inp{e_i}{e_j} = E_{ij}$
    $$
        E = \begin{pmatrix}
            E_{11} & \tikzmark{left}E_{12}\tikzmark{right} & \cdots & E_{1n} \\
            \DrawBox[thick]
            \tikzmark{left}E_{21}\tikzmark{right} & E_{22} & \cdots & E_{2n} \\
                \vdots & \vdots & \ddots & \vdots \\
                E_{n1} & E_{n2} & \cdots & E_{nn}
            \end{pmatrix}
            \DrawBox[thick]
    $$
\end{frame}
\begin{frame}{Inner Product as a Matrix}
    The Inner Product can be represented by a matrix very similarly to the Lie Bracket.
    \\$\inp{e_i}{e_j} = E_{ij}$
    $$
        E = \begin{pmatrix}
            E_{11} & E_{12} & \cdots & E_{1n} \\
            E_{12} & E_{22} & \cdots & E_{2n} \\
                \vdots & \vdots & \ddots & \vdots \\
                E_{1n} & E_{2n} & \cdots & E_{nn}
            \end{pmatrix}
    $$
\end{frame}
\begin{frame}{Inner Product as a Matrix}
    The Inner Product can be represented by a matrix very similarly to the Lie Bracket.
    \\$\inp{e_i}{e_j} = E_{ij}$
    $$
        E = \begin{pmatrix}
            E_{11} & E_{12} & \cdots & E_{1n} \\
            E_{12} & E_{22} & \cdots & E_{2n} \\
                \vdots & \vdots & \ddots & \vdots \\
                E_{1n} & E_{2n} & \cdots & E_{nn}
            \end{pmatrix}
    $$
    \pause
    Recall also that the inner product has additional properties.
    \pause
    \\Postive Definite: $\inp{x}{x} \ge 0 \ \text{and} \ \inp{x}{x}=0 \implies x=\zvec$
    \\And $\zvec$ canot be a basis vector because it is linearly dependent with all vectors.
    \\So $\blu{e_{ii}>0}$
\end{frame}
\begin{frame}{Inner Product as a Matrix}
    The Inner Product can be represented by a matrix very similarly to the Lie Bracket.
    \\$\inp{e_i}{e_j} = E_{ij}$
    $$
        E = \begin{pmatrix}
            \blu{E_{11}} & E_{12} & \cdots & E_{1n} \\
            E_{12} & \blu{E_{22}} & \cdots & E_{2n} \\
            \vdots & \vdots & \blu{\ddots} & \vdots \\
            E_{1n} & E_{2n} & \cdots & \blu{E_{nn}}
            \end{pmatrix}
    $$
    Recall also that the inner product has additional properties.
    \\Postive Definite: $\inp{x}{x} \ge 0 \ \text{and} \ \inp{x}{x}=0 \implies x=\zvec$
    \\And $\zvec$ canot be a basis vector because it is linearly dependent with all vectors.
    \\So $\blu{e_{ii}>0}$
\end{frame}
\begin{frame}{Dot Product on $\RR^3$}
    The Dot Product can be represented in this way.
    \pause
    \\$\inp{e_i}{e_j} = 0$
    \pause
    \\$\inp{e_i}{e_i} = 1$
    \pause
    $$
        E = \begin{pmatrix}
                1 & 0 & 0\\
                0 & 1 & 0\\
                0 & 0 & 1\\
            \end{pmatrix}
    $$
\end{frame}

\begin{frame}{$j$-map as a matrix}
Recall that for any $z \in \laz$, the map $j_z :\lav \to \lav$ is defined by
the identity
\begin{equation}
\inp{y}{j_z(x)}_\lav = \inp{z}{\br{x}{y}}_\laz.\label{jm}
\end{equation}
We wish to find a matrix $J_z$ that represents the map $j_z$. Since $j_z$ maps
$\lav$ to itself, $J_z$ should be a matrix in $\R^{n\times n}$. Rewriting
\eqref{jm} in matrix form, we get
    \begin{align*}
   \inp{y}{j_z(x)}_\lav & = \inp{z}{\br{x}{y}}_\laz \\
    y^T E (J_z x) 		& =  z^\t(x^TLy) \\
    y^T (EJ_z) x		& =  z^\t(-y^TLx) \\
                        & =  - z^\t(y^TLx) \\
                        & =  -y^T (z^\t L)x \\
                        & =  y^T (z^\t (-L)) x \\
                        & =  y^T (z^\t L^\T) x.
    \end{align*}
\end{frame}
\begin{frame}
The equation $y^T(EJ_z)x = y^T(z^\t L^\T) x$ must hold for arbitrary vectors
$x$ and $y$ in $\lav$, implying that
$$
EJ_z = z^\t L^\T.
$$
Since $\det(E) \neq 0$, it is invertible and
\begin{equation}
J_z = E^{-1}(z^\t L^\T).
\end{equation}
This gives us a method to compute the $j$-map for any $z$ using only matrix
computations. Furthermore, if $z = z_k$ is a basis vector of $\laz$, then
$z_k^\t L^\T$ is just $(L^k)^\T$. Thus, for basis vectors $z_k$, the map
$j_{z_k}$ is represented by the matrix
$$
J_{z_k} = E^{-1}(L^k)^\T.
$$
\end{frame}
\begin{frame}
Since the $j$-maps are linear, knowing how they act on a basis of $\laz$ is
good enough. Indeed, we should be able to define a stack of $j$-maps by
\begin{equation}
J = E^{-1}L^\T.
\end{equation}
The $j$-map $j_z$ \emph{should} then be $J_z = z^\t J$.
\end{frame}

\begin{frame}{What we did}
    Using the programming language, sage, we created a program that would
    compute the $j$-maps of an arbitrary Lie algebra.

\end{frame}

\section*{}
\subsection*{}

\begin{frame}{TikZ}
\smallskip

\smallskip

\end{frame}

\section*{}
\subsection*{}

\begin{frame}{Acknowledgements}
    \begin{figure}
    \centering
    \includegraphics[width=30mm]{lStack.tikz} %input file
    \label{fig:pic} %label name
    \caption{\footnotesize The L-Stack} %caption
    \end{figure}
    
        \begin{figure}
    \centering
    \includegraphics[width=30mm]{lStackNOTOP.tikz} %input file
    \label{fig:pic} %label name
    \caption{\footnotesize The L-Stack} %caption
    \end{figure}

\end{frame}

\end{document}
